\documentclass[10pt,a4paper]{amsart}
\usepackage{amsmath,amssymb,amsthm,graphicx}
\usepackage{verbatim}
\newenvironment{reference}{\comment}{\endcomment}
%\newenvironment{reference}{}{}
\newenvironment{slogan}{\comment}{\endcomment}
\newenvironment{history}{\comment}{\endcomment}

% For commutative diagrams we use Xy-pic
\usepackage[all]{xy}

% We use 2cell for 2-commutative diagrams.
\xyoption{2cell}
% Theorem environments.
%
\theoremstyle{plain}
\newtheorem{theorem}[subsection]{Theorem}
\newtheorem{proposition}[subsection]{Proposition}
\newtheorem{lemma}[subsection]{Lemma}

\theoremstyle{definition}
\newtheorem{definition}[subsection]{Definition}
\newtheorem{example}[subsection]{Example}
\newtheorem{exercise}[subsection]{Exercise}
\newtheorem{situation}[subsection]{Situation}

\theoremstyle{remark}
\newtheorem{remark}[subsection]{Remark}
\newtheorem{remarks}[subsection]{Remarks}

\numberwithin{equation}{subsection}

% Macros
%
\def\lim{\mathop{\mathrm{lim}}\nolimits}
\def\colim{\mathop{\mathrm{colim}}\nolimits}
\def\Spec{\mathop{\mathrm{Spec}}}
\def\Hom{\mathop{\mathrm{Hom}}\nolimits}
\def\Ext{\mathop{\mathrm{Ext}}\nolimits}
\def\SheafHom{\mathop{\mathcal{H}\!\mathit{om}}\nolimits}
\def\SheafExt{\mathop{\mathcal{E}\!\mathit{xt}}\nolimits}
\def\Sch{\mathit{Sch}}
\def\Mor{\mathop{\mathrm{Mor}}\nolimits}
\def\Ob{\mathop{\mathrm{Ob}}\nolimits}
\def\Sh{\mathop{\mathit{Sh}}\nolimits}
\def\NL{\mathop{N\!L}\nolimits}
\def\CH{\mathop{\mathrm{CH}}\nolimits}
\def\proetale{{pro\text{-}\acute{e}tale}}
\def\etale{{\acute{e}tale}}
\def\QCoh{\mathit{QCoh}}
\def\Ker{\mathop{\mathrm{Ker}}}
\def\Im{\mathop{\mathrm{Im}}}
\def\Coker{\mathop{\mathrm{Coker}}}
\def\Coim{\mathop{\mathrm{Coim}}}

% Boxtimes
%
\DeclareMathSymbol{\boxtimes}{\mathbin}{AMSa}{"02}

%
% Macros for moduli stacks/spaces
%
\def\QCohstack{\mathcal{QC}\!\mathit{oh}}
\def\Cohstack{\mathcal{C}\!\mathit{oh}}
\def\Spacesstack{\mathcal{S}\!\mathit{paces}}
\def\Quotfunctor{\mathrm{Quot}}
\def\Hilbfunctor{\mathrm{Hilb}}
\def\Curvesstack{\mathcal{C}\!\mathit{urves}}
\def\Polarizedstack{\mathcal{P}\!\mathit{olarized}}
\def\Complexesstack{\mathcal{C}\!\mathit{omplexes}}
% \Pic is the operator that assigns to X its picard group, usage \Pic(X)
% \Picardstack_{X/B} denotes the Picard stack of X over B
% \Picardfunctor_{X/B} denotes the Picard functor of X over B
\def\Pic{\mathop{\mathrm{Pic}}\nolimits}
\def\Picardstack{\mathcal{P}\!\mathit{ic}}
\def\Picardfunctor{\mathrm{Pic}}
\def\Deformationcategory{\mathcal{D}\!\mathit{ef}}
\usepackage{cleveref}
\crefname{originalsection}{original section}{original sections}
\Crefname{originalsection}{Original section}{Original sections}
\makeatletter
\newlabel{section-nonfree}{{065J}{source}{}{sourcepointer-1}{}}
\newlabel{section-nonfree@cref}{{[originalsection][0][]065J}{[1][0][]source}{}{}{}}
\makeatother
\title{A domain with a nonzero ideal principal at every prime but not invertible}
\author{The Stacks Project Authors}
\date{stacks-project:examples.tex:local-snapshot-20261012; GNU FDL 1.2 or later}
\begin{document}
\maketitle
\noindent\textit{Editorial note.} All delivered mathematical text is the author's own, copied from the pinned Stacks Project chapter file examples.tex; no statement or proof step was rewritten, reordered or removed.\par\medskip
\noindent\textit{Editorial note.} The counted proof is the author's complete argument: The terminology paragraph on 'locally principal' (2560-2571) is the author's terminological remark and is not needed for the argument; the counted chain uses only the two spans.\par\medskip
\noindent\textit{Editorial note.} The label lemma-locally-principal-not-invertible names the counted statement, and the proof environment immediately following it in the source is delivered verbatim as its proof.\par\medskip
\noindent\textit{Editorial note.} The counted statement is carried by the author's own lemma environment in the chapter source.\par\medskip
\noindent\textit{Allowed standard prerequisites (editorial).}\begin{itemize}
\item Context span 2536-2541 (section-nonfree): R = Q[x] and M = sum over n of (1/(x-n)) R as a submodule of the fraction field; M is not finitely generated, but for every prime p of R we have M\_p = R\_p as an R\_p-module (the author's stalk computation stated in the section)
\item Context span 2574-2597: A = Sym*\_R(M) (an example of a non-Noetherian domain with Noetherian local rings, author's footnote) and I = MA = the direct sum of the Sym\textasciicircum{}d\_R(M) for d >= 1; M not finitely generated implies I is not finitely generated as an ideal, and an invertible module is finitely generated, so I is not invertible; for every prime p of R, M\_p = R\_p gives A\_p = Sym*\_\{R\_p\}(R\_p) = R\_p[T] as a graded R\_p-algebra, so I\_p is generated by the nonzerodivisor T; hence for every prime q of A the ideal I\_q is generated by a single element
\end{itemize}\medskip
\ifdefined\hypertarget\hypertarget{sourcepointer-1}{}\fi
\noindent\textit{Source reference.} Original section 065J: sources/stacks-examples/pinned-original-source.tex, source lines 2533--2533. Prerequisite pointer invoked at 2623; it names a standard fact of the same chapter that the counted argument uses, and every step of the counted proof built on it is retained verbatim in the delivered body.\par\medskip
Let $R = \mathbf{Q}[x]$. Set
$M = \sum_{n \in \mathbf{N}} \frac{1}{x - n}R$ as a submodule of
the fraction field of $R$. Then $M$ is not finitely generated, but
for every prime $\mathfrak p$ of $R$ we have
$M_{\mathfrak p} \cong R_{\mathfrak p}$ as an
$R_{\mathfrak p}$-module.

\medskip\noindent
An example of a similar flavor is $R = \mathbf{Z}$ and
$M = \sum_{p \text{ prime}} \frac{1}{p} \mathbf{Z} \subset \mathbf{Q}$,
which equals the set of fractions $\frac{a}{b}$ with $b$ nonzero and
squarefree.

Let $R = \mathbf{Q}[x]$ and let $M = \sum \frac{1}{x - n}R$
be the module constructed in Section \ref{section-nonfree}.
Consider the ring\footnote{The ring $A$ is an example
of a non-Noetherian domain whose local rings are Noetherian.}
$$
A = \text{Sym}^*_R(M)
$$
and the ideal $I = M A = \bigoplus_{d \geq 1} \text{Sym}^d_R(M)$.
Since $M$ is not finitely generated as an $R$-module we see that
$I$ cannot be generated by finitely many elements as an ideal in $A$.
Since an invertible module is finitely generated, this means that
$I$ is not invertible.
On the other hand, let $\mathfrak p \subset R$ be a prime ideal.
By construction $M_\mathfrak p \cong R_\mathfrak p$. Hence
$$
A_\mathfrak p = \text{Sym}^*_{R_\mathfrak p}(M_\mathfrak p) \cong
\text{Sym}^*_{R_\mathfrak p}(R_\mathfrak p) =
R_\mathfrak p[T]
$$
as a graded $R_\mathfrak p$-algebra. It follows that
$I_\mathfrak p \subset A_\mathfrak p$ is generated by
the nonzerodivisor $T$.
Thus certainly for any prime ideal $\mathfrak q \subset A$
we see that $I_\mathfrak q$ is generated by a single element.

\begin{lemma}
\label{lemma-locally-principal-not-invertible}
There exists a domain $A$ and a nonzero ideal $I \subset A$
such that $I_\mathfrak q \subset A_\mathfrak q$ is a principal
ideal for all primes $\mathfrak q \subset A$ but $I$ is not an invertible
$A$-module.
\end{lemma}

\begin{proof}
See discussion above.
\end{proof}

\end{document}
